\section{Non-linearity and feature scale are indistinguishable}
\label{sec:confound}

Take equal linear effects but
skewed feature \emph{distributions}:
\[
  a = 1 + U_a, \quad b = (1 + U_b)^2, \quad c = (1 + U_c)^3, \quad y = a + b +
  c,
  \quad n = 10000,
\]
where $U_a, U_b, U_c$ are \emph{the same} draws used in
Section~\ref{sec:nonlinear}. Every feature enters with coefficient 1, and OLS
recovers $1.00, 1.00, 1.00$: nothing about the \emph{effects} is non-linear. Yet
the forest reports \textbf{1751 / 8166 / 39632}, a $1:4.7:22.6$ spread
--- exactly the same as the non-linear model of
Section~\ref{sec:nonlinear} (Figure~\ref{fig:confound}).

The agreement is exact because the two datasets present the forest with the
same information: the same response vector, and each feature in the same
order. The following proposition makes this precise.

\paragraph{Proposition.} Fix draws $U_a, U_b, U_c$ and compare two worlds:
\emph{(A)} features $(1{+}U_a,\; 1{+}U_b,\; 1{+}U_c)$ with response
$y = a + b^2 + c^3$, and \emph{(B)} features
$(1{+}U_a,\; (1{+}U_b)^2,\; (1{+}U_c)^3)$ with response $y = a + b + c$. Then
the
response vector is the same in both worlds, and each world-B feature is a
strictly increasing transform of its world-A counterpart. If (i) bootstrap
indices and candidate-feature draws are held fixed across the two fits, and
(ii) split selection depends on a feature only through the \emph{ordering} of
its values --- as it does for an exhaustive search over thresholds between
consecutive order statistics, \code{ranger}'s default
\code{splitrule = "variance"} --- then the two forests are identical tree for
tree as partitions of the in-bag sample, and every importance that is a
functional of those partitions (MDI,
AIR) coincides exactly.

Equal variance contribution is a by-product of the same construction; the
argument does not use it.

Both hypotheses can fail in practice. Hypothesis (i) fails in any script that
fits the two worlds one after the other without reseeding, and the cross-world
variation that follows is seed noise: world B under any RNG state gives the same
trees as world A under that state, so an unseeded cross-world pair is
distributed exactly as refitting a single world under two seeds. A sweep over
seeds $1$--$10$ confirms the corollary in \code{ranger}'s implementation: the
identity is exact at every seed, and across all $\binom{10}{2} = 45$ pairs of
distinct seeds the cross-world gap equals the same-world gap.
Hypothesis (ii) concerns the split search: a rule that sees a feature only
through the ordering of its values discards exactly the distribution-shape
information that would tell the two worlds apart.
\code{splitrule = "extratrees"} keeps a little of it, drawing thresholds
uniformly from the feature's value \emph{range} instead of scanning order
statistics. Squaring stretches that range, so the same draw lands elsewhere in
the data and the two worlds separate: $1:4.9:23.7$ against $1:5.0:24.6$.
The same thing happens to predictions. \citet{galili2016} show that a tree's
predictions for new observations are not invariant to monotone transforms of its
features, for the same reason: a threshold placed between two observed values
sits at their midpoint, and stretching the scale moves that midpoint even though
the training points either side of it stay put.

This is \textbf{not specific to impurity}, and debiasing does not touch it.
Bias-corrected impurity (AIR), permutation importance, and global mean
$|\text{SHAP}|$ all give the same ratios in both worlds:

\begin{center}
\begin{tabular}{lcc}
\toprule
measure & ratio & cross-world agreement \\
\midrule
impurity (MDI)            & $1:4.7:22.6$   & exactly $0$ \\
impurity, debiased (AIR)  & $1:8.8:47.4$   & exactly $0$ \\
permutation               & $1:9.2:50.3$   & ${\sim}10^{-4}$ \\
mean $|\text{SHAP}|$      & $1:3.06:6.98$  & $1.3 \times 10^{-4}$ \\
\midrule
standardised OLS          & $1:2.98:6.88$ (A), $1:2.99:6.98$ (B) & $1.4 \times
10^{-2}$ \\
\bottomrule
\end{tabular}
\end{center}

The first four follow from the proposition, in descending strictness. MDI is a
functional of the fitted trees alone, so it agrees identically. AIR agrees
identically for the same reason, one step removed: the correction subtracts the
impurity a permuted copy of each feature would earn, and under hypothesis (i)
those permutation draws are fixed too, while permuting commutes with monotone
transforms --- the permuted copy of $(1+U_b)^2$ is the square of the permuted
copy of $1+U_b$ --- so the correction is a functional of the same trees and the
same orderings. Debiasing moves the ratio from plain MDI's compressed
$1:4.7:22.6$ toward the variance-contribution ratio $1:9.07:48.96$ (a single
draw; Section~\ref{sec:linear} warns against reading any single-seed AIR
ratio closely and that applies here too), and moves the cross-world agreement
not at
all. Permutation and
path-dependent TreeSHAP \citep{lundberg2020} also touch feature values at the
thresholds --- one
through the values it shuffles, the other through node weights it reconstructs
from the data, which differ by at most two observations on $0.3\%$ of nodes ---
and agree to four significant figures rather than exactly.

Standardised OLS is the control. It never sees the forest, so the proposition
gives it nothing, and the two worlds are genuinely different regressions:
projecting $y$ onto $1+U$ is not projecting it onto $(1+U)^2$. The gap between
the two has a closed form. With independent features, the standardised
coefficient of $b$ is $\mathrm{sd}(b^2)/\mathrm{sd}(y)$ in world B and
$\mathrm{Cov}(b^2, b)/(\mathrm{sd}(b)\,\mathrm{sd}(y)) =
\mathrm{Corr}(b, b^2)\,\mathrm{sd}(b^2)/\mathrm{sd}(y)$ in world A, so the two
worlds differ by exactly the factor $\mathrm{Corr}(b, b^2)$: the fraction of the
squared feature that a linear projection captures. For $U(1,2)$ these
correlations are $0.9963$ for the square and $0.9861$ for the cube, and the
table reproduces them: $2.98/2.99 = 0.997$ and $6.88/6.98 = 0.986$. The
discrepancy is the population correlation gap itself, two orders of magnitude
larger than anything the proposition-pinned measures show. That a standardised
coefficient ties a variable's apparent importance to the spread of its
distribution is a long-standing complaint in the regression
literature \citep{bring1994}; the proposition is the tree-based version, with
effect shape and feature spread made exactly indistinguishable.

The spread of agreement is what identifies the mechanism. Equal variance
contribution alone would predict uniform approximate agreement, and cannot
produce an exact zero. Nor is being a variance the operative property: MDI and
permutation are denominated in variance, but mean $|\text{SHAP}|$ tracks the
\emph{square root} of the variance contribution ($\sqrt{1:9.07:48.96} =
1:3.01:7.00$), one power lower, and is confounded just the same. What the four
share, and what none of them escapes, is being a functional of the fitted model
and of the feature ordering, and the proposition makes both identical across the
two worlds. Any marginal importance inherits this,
whatever its units. The remedy is a shape-aware view: a partial-dependence plot,
or a SHAP \emph{dependence} plot, which plots each observation's SHAP value
against its feature value, and the feature values do differ between the two
worlds.

\begin{figure}[htb]
\centering
\includegraphics[width=0.85\textwidth]{fig3_confound.png}
\caption{Provably identical importances from two different worlds. Left:
genuinely
non-linear effects. Right: linear effects on skewed features. No marginal
importance measure can distinguish them.}
\label{fig:confound}
\end{figure}

\paragraph{Practical implications.} A tall importance bar can mean a strong or
non-linear effect, an interaction (Section~\ref{sec:interactions}), or merely a
heavy-tailed feature, and no marginal importance can say which --- not plain
or debiased impurity, not permutation, not mean $|\text{SHAP}|$. Switching
measure does not help; a shape-aware view does. The
partial-dependence plot \citep{friedman2001} of the skewed-feature world is a
straight line while the non-linear world's curves, at identical importances;
a SHAP dependence plot separates them the same way. The PDP averages the model
over the other features' marginal distribution, which is exact for
independent features as here; under dependence it evaluates the model on
combinations that do not occur in the data and can show its extrapolation
rather than an effect \citep{hooker2021}. There, accumulated local effects
(ALE) \citep{apley2020}, which
accumulates effects conditional on the neighbouring feature
distribution, replaces it. A SHAP dependence plot does not share this
exposure when the SHAP values are computed with path-dependent TreeSHAP, as
here: it is drawn only at observed data points, and path-dependent TreeSHAP
follows the training data down each tree rather than evaluating the model on
combinations that never occur.
